\chapter{Conclusiones y Recomendaciones}
\label{chap:conclusiones}

Este capítulo sostiene que el trabajo cumple lo que puede demostrar con evidencia verificable a la fecha de corte ---7 de noviembre de 2026--- y declara con la misma precisión lo que depende del levantamiento programado entre noviembre de 2026 y febrero de 2027. Por ello, cada conclusión responde a un objetivo específico, distingue entre resultados verificados y resultados por obtener, y no extiende sus afirmaciones más allá de lo que miden los instrumentos.

\section{Conclusiones}

\textbf{Respecto del objetivo específico 1 (requisitos e indicadores).} Se determinaron los requisitos de la LOPDP y de su Reglamento aplicables al tratamiento de datos de estudiantes y se tradujeron en treinta y un indicadores verificables: dieciséis documentales, observables en la información pública de las universidades, y quince de percepción, observables en la experiencia de los estudiantes (Tablas~\ref{tab:indicadores_documentales} y~\ref{tab:indicadores_percepcion}). Cada indicador quedó vinculado con su artículo, verificado en el texto oficial de la ley, y cada pregunta del cuestionario con su indicador (Tabla~\ref{tab:operacionalizacion}). El análisis mostró, además, que el núcleo verificable sin acceso interno es el derecho a la información del artículo~12, con sus diecisiete elementos, el consentimiento del artículo~8 y los derechos de los artículos~13 a 17, y que siete de esos elementos tienen equivalente en ambas fuentes, lo que hace posible medir la brecha entre lo publicado y lo percibido. Las obligaciones que solo pueden verificarse desde dentro ---registro de actividades, análisis de riesgos, evaluación de impacto, contratos y medidas técnicas--- se trasladaron a la metodología en lugar de suponerse cumplidas o incumplidas.

\textbf{Respecto del objetivo específico 2 (información pública de las universidades).} Se definió una población censal de 64 universidades y escuelas politécnicas (36 públicas, 20 particulares autofinanciadas y 8 cofinanciadas) según el CES, y se construyeron la ficha de verificación (Anexo~\ref{anx:ficha_documental}), el procedimiento de búsqueda y registro con evidencia fechada, el índice de transparencia documental (Ecuación~\ref{eq:itd}) y la plantilla de registro, cuyas 1\,410 fórmulas se evalúan sin errores y reproducen los cálculos independientes con datos sintéticos de prueba. El grado de transparencia de las universidades no puede afirmarse todavía, porque el análisis documental se ejecuta entre el 16 de noviembre y el 24 de diciembre de 2026; las Tablas~\ref{tab:res_documental} y~\ref{tab:itd_sostenimiento} recibirán sus resultados sin cambios de estructura.

\textbf{Respecto del objetivo específico 3 (encuesta a estudiantes).} Se diseñó un cuestionario anónimo de 46 preguntas cerradas ---dieciséis de hechos y veintitrés ítems de Likert en cuatro dimensiones, con cinco ítems inversos--- y un bloque de acceso que excluye a menores de edad y a quienes no consienten; se calculó un tamaño de muestra de referencia de 384 estudiantes para una población de 841\,404 matriculados (Ecuación~\ref{eq:muestra}) y se declaró que el muestreo es no probabilístico, con la consecuencia de que sus resultados describen a la muestra y no a toda la población. Se definieron los criterios de validación ---V de Aiken mayor o igual que 0,80 con tres a cinco jueces, prueba piloto con unos 30 estudiantes, alfa de Cronbach y omega de McDonald mayores o iguales que 0,70 por dimensión y KR-20 para la escala dicotómica---, tres hipótesis con sus hipótesis nulas y el plan de análisis, implementado en scripts reproducibles de Python que superan doce pruebas con valores de referencia y se ejecutaron de principio a fin con datos simulados. La percepción de los estudiantes no se reporta todavía, porque la recolección se realiza entre el 7 de diciembre de 2026 y el 15 de enero de 2027.

\textbf{Respecto del objetivo específico 4 (evaluación de la metodología).} La metodología, producto del trabajo, quedó construida a partir de los objetivos 1 a 3 y especificada para su evaluación: medidas de la brecha ---cumplimiento documental, brecha normativa y brecha de percepción (Ecuación~\ref{eq:brecha})--- que alimentan la identificación de riesgos; seis fases equivalentes a ISO 31000/ISO/IEC 27005, ISO/IEC 27035-1:2023 y NIST SP 800-61r3; escalas ancladas en frecuencias anuales; una tabla de decisión que impide ocultar un daño muy grave para el titular tras una probabilidad baja, y catorce escenarios sectoriales, cada uno vinculado con el indicador que lo activa, su artículo, sus controles, su rol responsable, su regla de plazo y su indicador de seguimiento. Cinco de los diez indicadores de validación cumplen su meta con evidencia disponible ---requisitos atendidos, fórmulas sin error en los tres libros de cálculo, casos de prueba superados y trazabilidad completa de los escenarios---; los cinco restantes, que corresponden al panel Delphi y a los juicios de expertos, no tienen resultado a la fecha de corte. En consecuencia, el trabajo no afirma todavía cuál es la brecha ni que la metodología sea válida; afirma que ambas cosas se medirán con criterios explícitos fijados de antemano. La metodología no se aplicó en una institución real, porque no se tuvo acceso a datos internos de ninguna universidad; su validación descansa en escenarios derivados de los hallazgos, en el panel Delphi y en el juicio de expertos.

\textbf{Limitaciones.} Primera, \textbf{la metodología no se aplicó en una institución real} por falta de acceso a datos internos: el estudio no usa ningún dato interno de ninguna universidad, no se pronuncia sobre el cumplimiento de las obligaciones que solo pueden verificarse desde dentro y valida la metodología con escenarios derivados de los hallazgos y con expertos, no con su aplicación en un proceso institucional. Segunda, la encuesta es no probabilística y se difunde por canales propios, de modo que sus resultados no se generalizan a todos los estudiantes del país. Tercera, la percepción del titular puede diferir de la práctica de la institución, aunque esa diferencia es en sí misma un dato sobre la transparencia. Cuarta, los sitios web cambian, por lo que cada resultado documental vale para su fecha de consulta. Quinta, la cifra de matrícula procede de una publicación que cita al SIIES, porque la base oficial no pudo descargarse en la fecha de consulta; el tamaño de muestra, sin embargo, no cambia con cualquier población superior a 100\,000 estudiantes. Sexta, la normativa de la SPDP está en evolución, de modo que el trabajo refleja la vigente a la fecha de cierre.

\section{Recomendaciones}

Las recomendaciones se derivan de los hallazgos verificados hasta la fecha y se completarán con los resultados del levantamiento:

\begin{enumerate}
  \item A las universidades: publicar una política de protección de datos que cubra los diecisiete elementos del artículo~12 de la LOPDP, con un canal operativo para ejercer derechos, y comprobar su propia transparencia aplicando la ficha documental de este trabajo a su sitio web.
  \item A las universidades: incorporar en la fase~0 de su gestión de riesgos una encuesta anónima a sus propios estudiantes con el cuestionario propuesto, porque una información que el titular no percibe no cumple el principio de transparencia.
  \item A las universidades: adoptar la plantilla de matriz como registro de riesgos, empezando por los escenarios que el diagnóstico active, y aprobar un procedimiento de respuesta a vulneraciones con los canales vigentes (SISPDP, organismo de control y CSIRT) y una meta interna de 72 horas.
  \item A la autoridad de control y a los organismos del sistema de educación superior: publicar, con fecha de corte, el número de delegados registrados por sector y las vulneraciones notificadas por el sector educativo, que la ley ordena registrar y que hoy no son públicas.
  \item Al autor y a la tutora: verificar los requisitos éticos de la universidad para la encuesta antes de su difusión y obtener la aprobación del nuevo título y de los objetivos.
\end{enumerate}

\section{Trabajos Futuros}

Se identifican cuatro líneas de continuidad. Primera, repetir el censo documental cada año con la misma ficha para medir la evolución del cumplimiento. Segunda, aplicar la metodología en una universidad que autorice el acceso a su información interna, para validar en un caso real las obligaciones que este trabajo no pudo verificar. Tercera, ampliar la encuesta con un muestreo probabilístico en cooperación con las universidades o con la SENESCYT, lo que permitiría estimar proporciones nacionales con un margen de error conocido. Cuarta, extender los indicadores a los institutos técnicos y tecnológicos y a la educación básica y el bachillerato, donde los titulares son mayoritariamente menores de edad.

En síntesis, el trabajo entrega un sistema de medición del cumplimiento de la LOPDP con fuentes legítimas, una metodología derivada de él y las herramientas para aplicarlos; lo que queda por demostrar depende de datos que se recogerán en el periodo previsto, y el documento los identifica uno por uno en lugar de suplirlos con supuestos.
